% ============================================================
% Come si usano, argomento 25: informazione ed entropia
% ============================================================
\argomento{Informazione ed entropia}{Formule a p.~\pageref*{f:informazione}}

\begin{scelta}{Come scegliere la formula}
\renewcommand{\arraystretch}{1.25}
\begin{tabular}{@{}P{0.55\linewidth}P{0.39\linewidth}@{}}
  \toprule
  \textbf{Se il testo\ldots} & \textbf{si usa\ldots} \\
  \midrule
  chiede quanta informazione dà un evento & \vv{informazione} \\
  chiede l'informazione media o l'incertezza di una sorgente, o i bit per simbolo necessari & \vv{entropia}; con due
    esiti \vv{entropiabinaria} \\
  chiede l'entropia più grande possibile con un certo alfabeto & \vv{entropiamax} \\
  confronta $E[g(X)]$ con $g(E[X])$ & \vv{jensen}, con \vv{convessa} \\
  \bottomrule
\end{tabular}

Le calcolatrici danno il logaritmo naturale: per i bit si usa $\log_2 x = \frac{\ln x}{\ln 2}$ (\vv{logbase2}). Se le
probabilità sono potenze di $\frac{1}{2}$, i logaritmi sono interi e i conti si fanno a mente.
\end{scelta}

\scheda{informazione}{
  \usoquando{Si chiede quanta informazione porta la notizia che un evento è avvenuto.}
  \usocome{Si calcola la probabilità dell'evento e si prende il logaritmo del reciproco, in base 2 per avere bit.}
  \usoesempio{Moneta onesta: 1 bit; faccia di un dado: $\log_2 6 \simeq 2{,}58$ bit; una carta da un mazzo di 52:
    $\log_2 52 \simeq 5{,}7$ bit; dieci teste di fila: $\log_2 2^{10} = 10$ bit, uno per lancio.}
}

\scheda{entropia}{
  \usoquando{Si chiede l'informazione media di una sorgente o di una variabile, o di quanti bit per simbolo c'è
    bisogno per rappresentarla.}
  \usocome{Per ogni valore si moltiplica la probabilità per l'informazione, e si somma.}
  \usoesempio{Sorgente $\left( \frac{1}{2}, \frac{1}{4}, \frac{1}{8}, \frac{1}{8} \right)$: $H = \frac{1}{2} \cdot 1
    + \frac{1}{4} \cdot 2 + \frac{1}{8} \cdot 3 + \frac{1}{8} \cdot 3 = 1{,}75$ bit, quanto la lunghezza media del
    codice \texttt{0}, \texttt{10}, \texttt{110}, \texttt{111} (\vv{codice}).}
}

\scheda{entropiabinaria}{
  \usoquando{L'esperimento ha due esiti: una moneta, una sorgente binaria, un bit che può essere sbagliato.}
  \usoesempio{Moneta truccata con $p = 0{,}11$: $H_2 \simeq 0{,}11 \cdot 3{,}18 + 0{,}89 \cdot 0{,}17 \simeq 0{,}5$ bit,
    quindi le sue sequenze di lanci si comprimono, in media, a metà.}
}

\scheda{jensen}{
  \usoquando{Si confronta la media di una funzione con la funzione della media, per esempio per dimostrare i limiti
    dell'entropia.}
  \usocome{Si controlla il segno della derivata seconda di $g$ sui valori di $X$ (\vv{convessa}).}
  \usoesempio{$g(x) = x^2$ è convessa: $E[X^2] \geq (E[X])^2$, cioè la varianza non è negativa. $X$ uniforme su $\{1,
    2, 3, 4\}$: $E[\log_2 X] \simeq 1{,}15 \leq \log_2 2{,}5 \simeq 1{,}32$, perché il logaritmo è concavo.}
}

\scheda{entropiamax}{
  \usoquando{Si chiede l'entropia più grande possibile con un certo numero di simboli, o se un codice a lunghezza fissa
    è già il migliore.}
  \usoesempio{Quattro simboli: al più $\log_2 4 = 2$ bit, con simboli equiprobabili. La sorgente $\left( \frac{1}{2},
    \frac{1}{4}, \frac{1}{8}, \frac{1}{8} \right)$ ha $1{,}75$ bit: il codice a lunghezza variabile sfrutta proprio la
    differenza.}
}
